\section{\ComputationCircuits{}}

We here suggest a quantum pendant to basis encodings (see Chapter~14 in \cite{goessmann_tensor-network_2025}), which has the decomposition by contraction property.

\begin{definition}[\ComputationCircuit{}]
    Given a boolean function $\exformula:\atomstates\rightarrow[2]$ the \computationCircuit{} is the unitary tensor
    \begin{align*}
        \qcbencodingofat{\exformula}{\cheadvariableof{\exformula}{\insymbol},\cheadvariableof{\exformula}{\outsymbol},\shortcatvariables}
        &\quad=
        \sum_{\shortcatindicesin\wcols\exformulaat{\shortcatindices}=1}
        \paulixat{\cheadvariableof{\exformula}{\insymbol},\cheadvariableof{\exformula}{\outsymbol}} \otimes
        \onehotmapofat{\shortcatindices}{\shortcatvariables} \\ % \otimes \onehotmapofat{\shortcatindices}{\ccatvariableof{[\atomorder]}{\outsymbol}} \\
        & \quad \quad \quad +
        \sum_{\shortcatindicesin\wcols\exformulaat{\shortcatindices}=0}
        \identityat{\cheadvariableof{\exformula}{\insymbol},\cheadvariableof{\exformula}{\outsymbol}} \otimes
        \onehotmapofat{\shortcatindices}{\shortcatvariables} \, . % \otimes \onehotmapofat{\shortcatindices}{\ccatvariableof{[\atomorder]}{\outsymbol}}
    \end{align*}
\end{definition}

%% WRONG
%Notice, that $\qcbencodingof{\lnot} = \mathrm{CNOT}$, which is obvious from $\onesat{\headvariableof{\insymbol},\headvariableof{\outsymbol}} - \identityat{\headvariableof{\insymbol},\headvariableof{\outsymbol}}$ being the Pauli-X gate (not to be confused with $X$ denoting distributed variables here).
%The \computationCircuit{} is therefore a generalized controlled $\mathrm{NOT}$ gate, where the control is by a boolean function.

\begin{example}[CNOT gate]\label{exa:cnotBenc}
    The CNOT gate is the controlled Pauli-X, that is
    \begin{align*}
        \mathrm{CNOT}\left[\cheadvariable{\insymbol},\cheadvariable{\outsymbol},\catvariable\right]
        = \paulixat{\cheadvariable{\insymbol},\cheadvariable{\outsymbol}} \otimes \tbasisat{\catvariable}
        + \identityat{\cheadvariable{\insymbol},\cheadvariable{\outsymbol}} \otimes \fbasisat{\catvariable} \, .
    \end{align*}
    For the identity function
    \begin{align*}
        \mathrm{Id}(x) = x \, .
    \end{align*}
    we have
    \begin{align*}
        \qcbencodingofat{\mathrm{Id}}{\cheadvariable{\insymbol},\cheadvariable{\outsymbol},\catvariable}
        = \mathrm{CNOT}\left[\cheadvariable{\insymbol},\cheadvariable{\outsymbol},\catvariable\right] \, .
    \end{align*}
    The $\mathrm{CNOT}$ is therefore the \computationCircuit{} to the identity function.
\end{example}

% Multiple output variables
Functions with multiple output variables, i.e. $\exformula:\atomstates\rightarrow\bigtimes_{\selindexin}[2]$, can be encoded image coordinate wise as a concatenation of the respective circuits.

\subsection{Relation with basis encodings}

\BasisEncodings{} are schemes to encode functions into directed boolean tensors, which enable reasoning about functions by tensor network contractions (see \cite{goessmann_tensor_2026}).
To connect with this formalism, let us now show the relation to the \computationCircuits{}.

\begin{remark}[Basis and Amplitude Encodings in Quantum Computation \cite{schuld_machine_2021} and \tnreason{}]
    % Compare with Quantum Computing
    Basis Encoding in Quantum Computation refers to the representation of classical $n$ bit strings by $n$ qubit basis states, and is called one-hot encoding in \tnreason{}.
    The Basis Encoding scheme in \tnreason{} goes beyond this scheme and also encodes subsets by sums of one-hot encodings to the members of the set.
    In this way, relations and functions are represented by boolean tensors and contraction of them is refered as \BasisCalculus{}.

    Amplitude Encoding in Quantum Computation refers to the storage of complex numbers in the amplitudes of quantum states.
    The pendant in \tnreason{} is the Coordinate Encoding scheme, where real numbers are stored in the coordinates of real-valued tensors.
    Compared to Amplitude Encoding, Coordinate Encoding does not have the normalization constraint of quantum states.
    The Amplitude Encoding of the square root of a probability distribution is sometimes called q-sample.
\end{remark}


\begin{lemma}
    \label{lem:bencComCircuit}
    For any boolean function we have (see \figref{fig:bencComCircuit})
    \begin{align*}
        \qcbencodingofat{\exformula}{\cheadvariableof{\exformula}{\insymbol},\cheadvariableof{\exformula}{\outsymbol},\shortcatvariables}
        = \contractionof{
            \bencodingofat{\exformula}{\headvariableof{\exformula},\shortcatvariables},
            \bencodingofat{\oplus}{\cheadvariableof{\exformula}{\outsymbol},\cheadvariableof{\exformula}{\insymbol},\headvariableof{\exformula}}
        }{\cheadvariableof{\exformula}{\insymbol},\cheadvariableof{\exformula}{\outsymbol},\shortcatvariables} \, .
    \end{align*}
\end{lemma}

% Sparse decomposition
In \cite{goessmann_tensor_2026} (and more detailed in \cite{goessmann_tensor-network_2025}) sparse decompositions of \basisEncodings{} have been studied.
Based on the equivalence of \computationCircuits{} and contracted \basisEncodings{}, stated by \lemref{lem:bencComCircuit} we transfer these decomposition schemes in the following.

\begin{figure}[t]
    \begin{center}
        \input{tikz_pics/computation-circuits/bencoding_computation_circuit}
    \end{center}
    \caption{Relation of computation circuit and basis encodings.
    The \computationCircuit{} of a function is the contraction of its \basisEncoding{} with the \basisEncoding{} of the mod2 sum $\oplus$.}
    \label{fig:bencComCircuit}
\end{figure}

% Good place?
\begin{example}[Multiple-controlled NOT]\label{exa:mcnotBenc}
    Let us consider the $\catorder$-ary controlled $\mathrm{NOT}$ gate
    \begin{align*}
        \mathrm{MCNOT}^{\catorder}\left[\cheadvariable{\insymbol},\cheadvariable{\outsymbol},\shortcatvariables\right]
        = \paulixat{\cheadvariable{\insymbol},\cheadvariable{\outsymbol}} \otimes \onehotmapofat{\onetuple{[\catorder]}}{\shortcatvariables}
        + \identityat{\cheadvariable{\insymbol},\cheadvariable{\outsymbol}} \otimes \left(\onesat{\shortcatvariables}-\onehotmapofat{\onetuple{[\catorder]}}{\shortcatvariables}\right) \, .
    \end{align*}
    This is the \computationCircuit{} of the $\catorder$-ary logical conjunction
    \begin{align*}
        \land^{\catorder} \defcols [2]^{\catorder} \rightarrow [2]
        \quad , \quad
         \land^{\catorder}(\shortcatindices)
        = \begin{cases}
              1 & \ifspace \shortcatindices = \onetuple{[\catorder]} \\
              0 & \text{else}
        \end{cases} \, .
    \end{align*}
    Special instances are the $\mathrm{CNOT}$ for $\catorder=1$ (see \exaref{exa:cnotBenc}) and the $\mathrm{Toffoli}$ gate for $\catorder=2$.
\end{example}

\subsection{Construction for mod2-basis+ CP decompositions - Exploiting \PolynomialSparsity{}}

\red{We use here the Algebraic Normal Form (ANF) of the booleans!}

\begin{lemma}
    \label{lem:comCircuitCommutative}
    The composition of two \computationCircuits{} to functions $\exfunction$ and $\secexfunction$ (see \figref{fig:mod2ComputationCircuit}), where the \computationCircuits{} share the target qubit, is the \computationCircuit{} of their sum mod2.
    In particular, \computationCircuits{} commute when they have the same target qubit.
\end{lemma}
\begin{proof}
    This follows from the representation of \computationCircuits{} by basis encodings of the corresponding function and the mod2 sum $\oplus$.
\end{proof}

\begin{figure}[t]
    \begin{center}
        \input{tikz_pics/computation-circuits/mod2_sum_computation_circuit}
    \end{center}
    \caption{Equivalence of the composition of \computationCircuits{} with the \computationCircuit{} of their mod2 sum (\lemref{lem:comCircuitCommutative}).}
    \label{fig:mod2ComputationCircuit}
\end{figure}

Concatenating two \computationCircuit{}s, which have the same head qubit, is the \computationCircuit{} of their mod2 sum.

Note that in general controlled unitaries on the same target qubit do not commute.

A basis+ elementary function can be encoded by a single controlled NOT operation with auxiliary $\catvariable$ qubits.

This motivates the mod2-basis+ CP decomposition of tensors, which is exactly the decomposition of boolean polynomials into monomials.
Each monomial is called a terms (products of x or (1-x) factors), and minterms in case that all variables appear.

\begin{definition}
    Given a boolean tensor $\hypercore$, a mod2-basis+ CP decomposition is a collection $\sliceset$ of tuples $(\variableset,\catindexof{\variableset})$ with such that for any $\shortcatindicesin$
    \begin{align*}
        \hypercoreat{\indexedshortcatvariables}
        = \bigoplus_{(\variableset,\catindexof{\variableset})\in\sliceset} \contractionof{\onehotmapofat{\catindexof{\variableset}}{\catvariableof{\variableset}}}{\indexedshortcatvariables} \, .
    \end{align*}
\end{definition}

Using that basis CP decompositions are a special case of basis+ CP decompositions, we get the following rank bound.

\begin{lemma}
    The mod2-basis+ CP rank is bounded by the basis CP rank.
\end{lemma}
\begin{proof}
    Use $\variableset=[\atomorder]$, and $\catindexof{\variableset}$ to each supported state.
    Then the mod2-sum is a usual sum and the basis CP decomposition is also a mod2-basis+ CP decomposition.
\end{proof}

This shows in particular, that any propositional formula can be represented by a mod2-basis+ CP decomposition.

\begin{lemma}
    The \computationCircuit{} to a boolean tensor $\hypercore$ with a mod2-basis+ CP decomposition $\sliceset$ obeys
    \begin{align*}
        &\qcbencodingofat{\hypercore}{
            \cheadvariableof{\hypercore}{\insymbol},
            \cheadvariableof{\hypercore}{\outsymbol},
            \shortcatvariables
        } \\
        & \quad =
        \breakablecontractionof{
            \{\identityat{\cheadvariableof{\hypercore}{\insymbol},\headvariableof{0}},\identityat{\cheadvariableof{\hypercore}{\outsymbol},\headvariableof{\cardof{\sliceset}-1}}\} \cup
            \{\qcbencodingofat{\onehotmapof{\catindexof{\variableset}}}{
                \headvariableof{\decindex}, \headvariableof{\decindex+1},
                \catvariableof{\variableset}} \wcols (\variableset,\catindexof{\variableset})\in\sliceset\}
        }{\cheadvariableof{\hypercore}{\insymbol},\cheadvariableof{\hypercore}{\outsymbol},\shortcatvariables}
    \end{align*}
    where $\decindex\in[\cardof{\sliceset}]$ enumerates the tuples in $\sliceset$.
\end{lemma}
\begin{proof}
    Can be seen by the basis encoding representation, and the accociativtiy of the sum mod2 ($\oplus$).
\end{proof}

Each \computationCircuit{} to each boolean monomial can be prepared by a multiple-controlled $\paulixsymbol$ gate and further pairs of $\paulixsymbol$ gates preparing the control state, see \figref{fig:qcbencodingPolynomial}. %where the control qubits are the affected variables and the target qubit is the value qubit.
When we sum monomials wrt modulus 2 calculus, then the preparation is a sequence of such circuits.
In such way, we can prepare the \computationCircuit{} to any propositional formula.
This encoding strategy exploits a modified (by mod2 calculus) \polynomialSparsity{}.

\begin{figure}
    \begin{center}
        \input{tikz_pics/computation-circuits/slice_computation_circuit}
    \end{center}
    \caption{
        Representation of an slice as a \red{state-controlled} $\mathrm{CNOT}$ gate.
        \red{Notation by white dots is also problematic: Nielsen Chuang uses them for $0$-state.}
%        Exploitation of \PolynomialSparsity{} in \computationCircuit{}s.
    %Here, we use the typical denotation of a muliple-controlled CNOT gate, which control symbols do not indicate Dirac tensors.
    }\label{fig:qcbencodingPolynomial}
\end{figure}

\subsection{Composition by Contraction - Exploiting \DecompositionSparsity{}}

The decomposition by contraction property of basis encodings is now a composition of circuits property, as stated in the next lemma.
Instead of hidden variables as in the \tnreason{} approach, we here exploit further working qubits.

\begin{lemma}
    \label{lem:comCirDecSpar}
    We have for functions $\exfunction:\atomstates\rightarrow\bigtimes_{\selindexin}[2]$, $\secexfunction:\bigtimes_{\selindexin}[2]\rightarrow\bigtimes_{s\in[r]}[2]$ (see \figref{fig:qcbencodingDecomposition})
    \begin{align*}
        &\qcbencodingofat{\secexfunction\circ\exfunction}{
            \cheadvariableof{\secexfunction\circ\exfunction}{\insymbol},
            \cheadvariableof{\secexfunction\circ\exfunction}{\outsymbol},
            \shortcatvariables}
        \otimes \fbasisat{\cheadvariableof{\exfunction}{\outsymbol}} \\
        &\quad= \breakablecontractionof{\onehotmapofat{0}{\cheadvariableof{\exfunction}{\insymbol}},
            \qcbencodingofat{\exfunction}{\cheadvariableof{\exfunction}{\insymbol},
                \cheadvariableof{\exfunction}{\midsymbol},
                \shortcatvariables},
            \qcbencodingofat{\secexfunction}{
                \headvariableof{\insymbol,\secexfunction\circ\exfunction},
                \cheadvariableof{\secexfunction\circ\exfunction}{\outsymbol},
                \cheadvariableof{\exfunction}{\midsymbol}},
            \qcbencodingofat{\exfunction}{\cheadvariableof{\exfunction}{\midsymbol},
                \cheadvariableof{\exfunction}{\outsymbol},
                \shortcatvariables}
        }{
            \cheadvariableof{\secexfunction\circ\exfunction}{\insymbol},\cheadvariableof{\secexfunction\circ\exfunction}{\outsymbol},\shortcatvariables,\cheadvariableof{\exfunction}{\outsymbol}
        } \, .
    \end{align*}
    Here each copy of $\headvariableof{\secexfunction\circ\exfunction}$ denotes a tuple of $r$ boolean variables and $\headvariableof{\exfunction}$ of $\seldim$ boolean variables.
\end{lemma}
\begin{proof}
    We show the claim based on basis calculus and sketch the steps in \figref{fig:proofDecomposition}.
    In the first equation we use \lemref{lem:bencComCircuit} and represent each \computationCircuit{} by a contraction of two basis encodings.
    In the second equation we use the identity
    \begin{align*}
        \contractionof{\bencodingofat{\oplus}{\cheadvariable{\outsymbol},\cheadvariable{\insymbol},\headvariableof{\exfunction}},
            \fbasisat{\cheadvariable{\insymbol}}}{\cheadvariable{\outsymbol},\headvariableof{\exfunction}}
        = \identityat{\cheadvariable{\outsymbol},\headvariableof{\exfunction}} \, .
    \end{align*}
    Now in the third equation we use that the contraction of basis encodings is the basis encoding of the composition function, and that $\exfunction\oplus\exfunction$ is the vanishing function, which basis encoding is a tensor product with $\fbasisat{\cheadvariableof{\exfunction}{\outsymbol}}$.
    We use we \lemref{lem:bencComCircuit} again in the fourth equation to arrive at the claim.
%    For any $\shortcatindices\in[2]^{\atomorder}$ we have that (see \figref{fig:proofDecomposition}):
%    \red{Uses the polynomial sparsity to uncompute: $\exfunction \oplus\exfunction = 0$.}
%    Can be seen by the basis encoding representation, when starting the hidden qubit with $\fbasis$, since
%    \begin{align*}
%        \contractionof{\bencodingofat{\oplus}{\headvariable,X_1},\fbasisat{X_0}}{Y,X_1}
%        = \identityat{Y,X_1} \, .
%    \end{align*}
%    The marginalization of the hidden target qubit $X_1$ for measurements does not change the distribution, since we have a deterministic dependence on the measured data register (need preparation by $\fbasis$ for that).
\end{proof}

When having a syntactical decomposition of a propositional formula, we can iteratively apply the \computationCircuit{} decomposition theorem and prepare each connective by a circuit.
We can decompose any propositional formula into logical connectives and prepare to each a \computationCircuit{} (exploiting further sparsity mechanisms).
This works, when the target qubit of one connective is used as a value qubit of another.

%% Closure of
Note that to exploit \decompositionSparsity{}, we introduce another auxiliary working qubit representing the satisfaction of the sub-function $\exfunction$.
In \lemref{lem:comCirDecSpar} the equality of the concatenated \computationCircuits{} holds after contraction of this qubit, see also \figref{fig:qcbencodingDecomposition}.
When performing a computational basis measurement on the state prepared by a concatenation, and omitting the qubit (either not measure it in the first place, or ignore its measurement), one effectively has closed the variable before measurement.
This holds true due to the deterministic dependence of the measurement outcome on the measurement of the distributed qubits $\shortcatvariables$, but is not true for more generic concatenations of controlled unitaries.
This entailment of the working qubit with the distributed qubits can be resolved by an un-computation circuit, which is the adjoint of the \computationCircuit{} for the working qubits (see \figref{fig:qcbencodingDecomposition}).
Note that while this is necessary in some applications, it is not necessary for the quantum rejection sampling presented here.

\begin{figure}
    \begin{center}
        \input{tikz_pics/computation-circuits/decomposition_sparsity}
    \end{center}
    \caption{
        Exploitation of \DecompositionSparsity{} in \computationCircuit{}s.
        The working qubit $\headvariableof{\exfunction}$ is initialized in the ground state $\fbasis$ and un-computed after its usage to prepare $\headvariableof{\secexfunction\circ\exfunction}$.
        It is then disentangled with the other qubits.
    }\label{fig:qcbencodingDecomposition}
\end{figure}

\begin{figure}
    \begin{center}
        \input{tikz_pics/computation-circuits/proof_decomposition_sparsity}
    \end{center}
    \caption{
        Proof of \lemref{lem:comCirDecSpar} based on basis calculus.
    }\label{fig:proofDecomposition}
\end{figure}

\begin{example}\label{exa:mcnotDecomposition}
    Let us reconsider the $\catorder$-ary controlled $\mathrm{NOT}$ gate from \exaref{exa:mcnotBenc}.
    The $\catorder$-ary logical conjunction is decomposed as a sequence of $2$-ary logical conjunctions.
    The \computationCircuit{} of each $2$-ary conjunction is a $\mathrm{Toffoli}$ gate applied to an input qubit and a worker qubit.
    Using \DecompositionSparsity{} we therefore find the standard decomposition of $\catorder$-ary controlled $\mathrm{NOT}$ gates into $\mathrm{Toffoli}$ gates.
\end{example}

\subsection{Preparation by fine and coarse structure}

Having a mod2-basis+ CP decomposition of rank $r$ to a connective, we need $r$ controlled NOT gates to prepare the basis encoding.
Given a syntactical decomposition of a boolean statistics, we prepare the basis encoding as a circuit with:
\begin{itemize}
    \item \textbf{Fine Structure:} Represent each logical connective based on its mod2-basis+ CP decomposition, as a concatenation of \computationCircuit{}s with the same variables.
    \item \textbf{Coarse Structure:} Arrange the logical connective representing circuits according to the syntactical hypergraph, where parent head variables appear as distributed variables at their children.
\end{itemize}

